\expandafter\gdef\csname SNside1\endcsname{\node[anchor=north west,inner sep=0pt,text width=48mm,font=\SideFont,align=left] at ([xshift=138mm,yshift=-513.45pt]current page.north west) {\textcolor{SNink}{\textbf{课程连接 · 驱动不是设备本身}}\par\vspace{1.5mm}OS中的驱动程序在CPU上执行，I/O接口是硬件，外部设备是数据来源或去向。DMA替换逐次数据搬运的执行主体，仍须驱动设置与服务程序收尾，软硬件边界要逐阶段标出。};}
\expandafter\gdef\csname SNside2\endcsname{\node[anchor=north west,inner sep=0pt,text width=48mm,font=\SideFont,align=left] at ([xshift=138mm,yshift=-475.52pt]current page.north west) {\textcolor{SNink}{\textbf{最小反例 · 24位不等于24种颜色}}\par\vspace{1.5mm}每像素24位时最多有 \(2^{24}\) 种编码，而不是24种。容量计算乘24得到位数，再除8变字节；若题给的是灰度级G，应先算编码位数 \(\lceil\log_2G\rceil\)。\par\smallskip\begin{tikzpicture}[x=1mm,y=1mm,>=Stealth,node distance=5mm]
\node[draw=SNink!50,fill=SNblue,rounded corners=1mm,text width=34mm,align=center,inner sep=1.2mm,font=\fontsize{7.5}{9}\selectfont] (m0) at (0,0) {像素24位};
\node[draw=SNink!50,fill=SNblue,rounded corners=1mm,text width=34mm,align=center,inner sep=1.2mm,font=\fontsize{7.5}{9}\selectfont] (m1) at (0,-11) {2的24次方};
\node[draw=SNink!50,fill=SNblue,rounded corners=1mm,text width=34mm,align=center,inner sep=1.2mm,font=\fontsize{7.5}{9}\selectfont] (m2) at (0,-22) {除8得字节};
\draw[->,SNink] (m0) -- (m1);
\draw[->,SNink] (m0) edge[bend left=40] (m2);
\end{tikzpicture}
\par{\fontsize{7.5}{9}\selectfont 位数决定编码数，容量还须把bit换成B}};}
\expandafter\gdef\csname SNside3\endcsname{\node[anchor=north west,inner sep=0pt,text width=48mm,font=\SideFont,align=left] at ([xshift=138mm,yshift=-95.00pt]current page.north west) {\textcolor{SNink}{\textbf{不变量 · 先分字符库与屏幕帧}}\par\vspace{1.5mm}字符点阵库按字符种类、字模行列及每点位数统计；帧缓冲按当前屏幕像素数统计。即使屏上同一字符出现多次，字库也不必重复保存同一字模，两个容量不能直接互换。};}
\expandafter\gdef\csname SNside4\endcsname{\node[anchor=north west,inner sep=0pt,text width=48mm,font=\SideFont,align=left] at ([xshift=138mm,yshift=-95.00pt]current page.north west) {\textcolor{SNink}{\textbf{不变量 · 最长间隔受缓冲期限限制}}\par\vspace{1.5mm}R B/s填满q B缓冲需要q/R秒；定时查询必须在覆盖前读走数据，还要计入实际检查与搬运延迟。平均每秒查够次数并不自动保证某个最长间隔不会超过期限。};}
\expandafter\gdef\csname SNside5\endcsname{\node[anchor=north west,inner sep=0pt,text width=48mm,font=\SideFont,align=left] at ([xshift=138mm,yshift=-95.00pt]current page.north west) {\textcolor{SNink}{\textbf{课程连接 · 模式切换不必换进程}}\par\vspace{1.5mm}设备中断可在内核态完成处理再返回原进程。保存必要硬件返回状态、软件保护现场以及调度换进程是不同层次，不能把每次中断都算成完整进程上下文切换。};}
\expandafter\gdef\csname SNside6\endcsname{\node[anchor=north west,inner sep=0pt,text width=48mm,font=\SideFont,align=left] at ([xshift=138mm,yshift=-95.00pt]current page.north west) {\textcolor{SNink}{\textbf{错项排除 · 块完成一次，单位传送多次}}\par\vspace{1.5mm}1024 B块、4 B接口缓冲可分256个数据单位搬运，而正常完成通知可仅一次。DMA减少CPU事件数，不要求把整块压成一个总线数据拍；突发是否合并事务还需题设。};}
\expandafter\gdef\csname SNside7\endcsname{\node[anchor=north west,inner sep=0pt,text width=48mm,font=\SideFont,align=left] at ([xshift=138mm,yshift=-95.00pt]current page.north west) {\textcolor{SNink}{\textbf{课程连接 · CPU预算与总线预算分列}}\par\vspace{1.5mm}每块软件预处理/后处理影响CPU比例，DMA实际访存影响总线和主存占用。CPU可能在其私有Cache中继续运算，所以总线忙碌比例不能无条件等同CPU损失比例。};}
\expandafter\gdef\csname SNside8\endcsname{\node[anchor=north west,inner sep=0pt,text width=48mm,font=\SideFont,align=left] at ([xshift=138mm,yshift=-95.00pt]current page.north west) {\textcolor{SNink}{\textbf{最小反例 · 平均低占用也会超期限}}\par\vspace{1.5mm}每秒总共只需1ms服务，看似占用很低；若某次数据必须在10微秒内取走，而CPU最长关中断100微秒，仍可能丢数据。平均预算与最坏响应延迟必须分别核对。};}
\expandafter\gdef\csname SNside9\endcsname{\node[anchor=north west,inner sep=0pt,text width=48mm,font=\SideFont,align=left] at ([xshift=138mm,yshift=-95.00pt]current page.north west) {\textcolor{SNink}{\textbf{不变量 · 每层恢复自己的屏蔽字}}\par\vspace{1.5mm}嵌套服务像调用栈逐层返回。C退出应先恢复B的屏蔽状态，B再恢复A；直接恢复最外层全开放状态可能允许原本不该打断B的请求，破坏题设处理关系。};}
\expandafter\gdef\csname SNside10\endcsname{\node[anchor=north west,inner sep=0pt,text width=48mm,font=\SideFont,align=left] at ([xshift=138mm,yshift=-95.00pt]current page.north west) {\textcolor{SNink}{\textbf{最小反例 · 延迟不是服务时间}}\par\vspace{1.5mm}请求产生后等50ns才开始，服务程序运行100ns，真正占CPU执行的是后者。若这50ns期间CPU执行别的任务，就不能把它加入该设备服务时间；最长允许延迟只是期限。};}
\AddToShipoutPictureFG{\begin{tikzpicture}[remember picture,overlay]\ifcsname SNside\thepage\endcsname\csname SNside\thepage\endcsname\fi\end{tikzpicture}}